%% kapa: exam=1 round=19 year=2008 period=2 form=A subject=economics subject_ko=경제원론 seq=10 number=10
%% source: Q-Net 감정평가사 기출문제 (cst00309) artlSeq=1001962 file=extracted/1008/1001962/2008년 제19회 감정평가사 1차시험 기출문제/2교시.hwp
%% conversion: kordoc 4.24.51 -> tools/gen1.py (auto); review=proofread-sonnet; flags=none
\begin{question}{10}
프리드만의 항상소득가설에 대한 설명 중 가장 적절한 것은?
\begin{choices}
  \item 현재소비는 현재소득에만 의존한다.
  \item 항상소득과 일시소득 사이에 높은 상관관계가 있다.
  \item 현재소득이 일시적으로 항상소득 이상으로 증가하면 평균소비성향은 일시적으로 하락한다.
  \item 항상소득가설은 유동성제약이 없는 소비자의 경우보다 유동성제약이 있는 소비자의 경우에 부합한다.
  \item 소득의 변동이 일시적인 요인에 의한 경우에 소비의 변동이 소득의 변동보다 커진다.
\end{choices}
\end{question}
\answer{3}
